\documentclass[11pt,a4paper]{article}

\usepackage[utf8]{inputenc}
\usepackage[T1]{fontenc}
\usepackage{lmodern}
\usepackage{microtype}

\usepackage{amsmath,amssymb}
\usepackage{booktabs}
\usepackage{enumitem}
\usepackage{hyperref}
\usepackage{xcolor}
\usepackage{listings}
\usepackage[margin=1in]{geometry}

\hypersetup{
  colorlinks=true,
  linkcolor=blue!60!black,
  urlcolor=blue!60!black,
  citecolor=blue!60!black,
}

\lstset{
  language=C++,
  basicstyle=\ttfamily\small,
  keywordstyle=\color{blue!70!black}\bfseries,
  commentstyle=\color{green!50!black}\itshape,
  stringstyle=\color{red!60!black},
  numbers=none,
  frame=single,
  framerule=0.4pt,
  rulecolor=\color{black!30},
  backgroundcolor=\color{black!3},
  breaklines=true,
  breakatwhitespace=true,
  tabsize=4,
  showstringspaces=false,
  columns=flexible,
  aboveskip=1em,
  belowskip=1em,
}

\newcommand{\code}[1]{\texttt{#1}}
\newcommand{\file}[1]{\texttt{#1}}

\title{Assessment of Fitting Routines\\[4pt]
       \large epicChargeSharing}
\author{Code Review Notes}
\date{February 2026}

\begin{document}
\maketitle

%----------------------------------------------------------------------
\section{What You Have}
%----------------------------------------------------------------------

Three fitting codepaths:

\begin{table}[h]
\centering
\begin{tabular}{@{}lll@{}}
\toprule
\textbf{Component} & \textbf{File} & \textbf{Method} \\
\midrule
Standalone (post-processing) & \file{src/GaussianFitter.cc}          & Reads ROOT file, adds fit branches in batch \\
EICrecon plugin (inline)     & \file{eicrecon/include/GaussianFit.hh} & Per-event during reconstruction \\
Uncertainty weighting        & \file{include/FitUncertainty.hh}       & 9 distance-weighted $\sigma$ models \\
\bottomrule
\end{tabular}
\end{table}

\noindent
Both 1D and 2D Gaussian fits use:
\begin{itemize}[nosep]
  \item \textbf{Minimizer}: Fumili2 (fast, least-squares specialised) with Migrad fallback.
  \item \textbf{Model}:
    $f(x) = A \exp\!\bigl(-\tfrac{1}{2}\bigl(\tfrac{x-\mu}{\sigma}\bigr)^{2}\bigr) + B$
    \;(1D) or the 2D equivalent with 6 parameters.
  \item \textbf{Seeding}: $Q_{\max}$ for $A$, $Q_{\min}$ for $B$,
    position-of-max for $\mu$, weighted variance for $\sigma$.
  \item \textbf{Bounds}: $\mu$ constrained to $\pm 1$ pitch from the
    centre pixel; $\sigma \in [\text{pixelSize},\; R \cdot \text{pitch}]$.
  \item \textbf{Fallback}: Charge-weighted centroid when the fit fails.
\end{itemize}

%----------------------------------------------------------------------
\section{What's Done Well}
%----------------------------------------------------------------------

\paragraph{1.\ Minimizer strategy is textbook-correct.}
Fumili2 is the right first choice for $\chi^{2}$ minimisation in ROOT
(it exploits the sum-of-squares structure), and the Migrad fallback with
elevated strategy handles the cases Fumili2 can't.  This is what
experienced ROOT users do.

\paragraph{2.\ Initial parameter estimates are good.}
You seed from the data itself ($Q_{\min}$/$Q_{\max}$, position of max,
weighted variance).  Good seeds are 90\% of the battle with Minuit---many
HEP fits fail because of bad initial values, not because of bad
minimisers.

\paragraph{3.\ Thread safety is handled correctly.}
\code{thread\_local} TF1/TF2 objects, \code{FlatVectorStore} to decouple
TTree I/O from parallel fitting, \code{ROOT::TThreadExecutor} for
dispatch.  This is non-trivial to get right with ROOT.

\paragraph{4.\ Centroid fallback is essential and present}
in both standalone and plugin.  A fit failure shouldn't mean no
reconstruction.

\paragraph{5.\ Parameter bounds prevent pathological fits.}
The $\mu$ range ($\pm 1$ pitch) and $\sigma$ bounds are physically
reasonable constraints.

%----------------------------------------------------------------------
\section{The Fundamental Issue: Fitting the Wrong Function}
%----------------------------------------------------------------------

This is the single most important point.  The charge sharing profile from
your LogA model is:
\begin{equation}
  F_i \;\propto\; \frac{\alpha_i}{\ln(d_i / d_0)}
  \label{eq:loga}
\end{equation}
This is \textbf{not a Gaussian in distance}.  It falls off as
$1/\ln r$---much slower than $e^{-r^{2}}$.  A Gaussian fit to this data:

\begin{itemize}[nosep]
  \item Has \textbf{systematic bias} on the $\mu$ parameter due to model
    mismatch, particularly for off-centre hits.
  \item Produces \textbf{physically meaningless} $\sigma$ and $\chi^{2}$
    values.
  \item Requires \textbf{4 parameters for 5 data points} (1D: $R=2$ gives
    a central slice of 5 pixels, $\text{NDF}=1$)---the fit is essentially
    unconstrained.
  \item Drives the complexity in \file{FitUncertainty.hh}---those 9
    distance-weighted $\sigma$ models are compensating for the wrong fit
    function.
\end{itemize}

\noindent
The Gaussian works as a \emph{reasonable interpolator} (it captures the
peak location), which is why you get usable resolution.  But you're
leaving accuracy on the table.

%----------------------------------------------------------------------
\section{Recommended Improvement: Fit Your Own Physics Model}
\label{sec:physfit}
%----------------------------------------------------------------------

The core idea is simple: instead of fitting an arbitrary bell-shaped
curve (the Gaussian) to the observed charges and extracting the peak
position, fit the \emph{actual charge-sharing model} and extract the
hit position directly.  This section walks through both approaches
step-by-step so the contrast is clear.

%----------------------------------------------------------------------
\subsection{What the Current Gaussian Fit Does (Step by Step)}
\label{sec:current-fit}
%----------------------------------------------------------------------

Consider a single event.  A particle hits the sensor at some unknown true
position $(x_{\text{true}}, y_{\text{true}})$.  The simulation (or the
real detector) produces an observed charge $Q_i$ on each pixel~$i$ in
the neighbourhood.  For the 1D row fit with $R=2$, you extract the
central row ($dj = 0$) giving 5 data points:

\medskip
\begin{center}
\begin{tabular}{@{}cccccc@{}}
\toprule
Pixel $i$ & $-2$ & $-1$ & $0$ & $+1$ & $+2$ \\
\midrule
Position $x_i$ & $x_{\text{px}} - 2p$ & $x_{\text{px}} - p$
  & $x_{\text{px}}$ & $x_{\text{px}} + p$ & $x_{\text{px}} + 2p$ \\
Observed charge & $Q_{-2}$ & $Q_{-1}$ & $Q_0$ & $Q_{+1}$ & $Q_{+2}$ \\
\bottomrule
\end{tabular}
\end{center}
\medskip

\noindent
where $p$ is the pixel pitch and $x_{\text{px}}$ is the centre pixel
position.  You then define a fit function with \textbf{4 free parameters}
$(A, \mu, \sigma, B)$:
\begin{equation}
  f(x;\; A, \mu, \sigma, B)
    \;=\; A \exp\!\Bigl(-\frac{1}{2}\Bigl(\frac{x - \mu}{\sigma}\Bigr)^{\!2}\Bigr) + B
  \label{eq:gauss-model}
\end{equation}
and minimise:
\begin{equation}
  \chi^{2}
    \;=\; \sum_{i=-R}^{R}
      \frac{\bigl(Q_i - f(x_i;\; A, \mu, \sigma, B)\bigr)^{2}}
           {\sigma_{\text{err},i}^{\,2}}
  \label{eq:gauss-chi2}
\end{equation}
Minuit2 adjusts all four parameters until $\chi^{2}$ is minimised.  The
reconstructed $x$-position is the fitted~$\mu$.

\medskip\noindent
\textbf{The problem.}\quad
The function~$f$ is a Gaussian.  But the charges $Q_i$ were \emph{not}
generated by a Gaussian process---they were generated by your LogA (or
LinA) charge-sharing model.  The Gaussian has no physical connection to
the data.  It is just a smooth bell curve being used as a generic
interpolator.  Because the shapes don't match:

\begin{itemize}[nosep]
  \item The minimiser must use 4 degrees of freedom ($A, \mu, \sigma, B$)
    to approximate a shape that isn't actually Gaussian, leaving only
    $\text{NDF} = 5 - 4 = 1$.
  \item The parameters $\sigma$ and $B$ absorb the model mismatch, biasing
    the value of~$\mu$.
  \item The $\chi^{2}$ value doesn't tell you whether the fit is good,
    because the model is wrong by construction.
  \item You need elaborate $\sigma_{\text{err},i}$ models
    (\file{FitUncertainty.hh}) to partially compensate for the mismatch
    by re-weighting which pixels the fitter trusts.
\end{itemize}

%----------------------------------------------------------------------
\subsection{What the Physics-Model Fit Does Instead}
\label{sec:phys-fit}
%----------------------------------------------------------------------

The key insight: you already have a function that predicts exactly what
$Q_i$ should be for any trial hit position.  That function is
\code{calcWeight()} in \file{ChargeSharingCore.hh}.  So instead of
fitting $f(x; A, \mu, \sigma, B)$, fit the \emph{actual physics} with
the hit position as the only free parameter.

\paragraph{Step 1: Define the forward model.}
For a trial hit position $x_0$ (the single fit parameter), compute the
predicted charge on each pixel.  Each pixel centre is at a known position
$c_i$.  The distance from the trial position to pixel~$i$ is:
\begin{equation}
  d_i(x_0) \;=\; |x_0 - c_i|
  \qquad\text{(1D row fit)}
  \label{eq:dist-1d}
\end{equation}
For the full 2D fit with trial position $(x_0, y_0)$:
\begin{equation}
  d_i(x_0, y_0) \;=\; \sqrt{(x_0 - c_{i,x})^{2} + (y_0 - c_{i,y})^{2}}
  \label{eq:dist-2d}
\end{equation}
Then compute the weight for each pixel using your existing model.  For
LogA:
\begin{equation}
  w_i(x_0) \;=\; \frac{\alpha_i\bigl(d_i(x_0)\bigr)}
                       {\ln\bigl(d_i(x_0) / d_0\bigr)}
  \label{eq:weight-loga}
\end{equation}
where $\alpha_i$ is the pad view angle (also a function of $d_i$), and
$d_0$ is the reference distance.  These are exactly the functions
\code{calcPadViewAngle()} and \code{calcWeightLogA()} that you already
have.

\paragraph{Step 2: Normalise to get predicted fractions.}
\begin{equation}
  F_i^{\,\text{pred}}(x_0) \;=\;
    \frac{w_i(x_0)}{\displaystyle\sum_{j} w_j(x_0)}
  \label{eq:fraction}
\end{equation}
This is the same normalisation your simulation already performs in
\code{calculateNeighborhood()}.

\paragraph{Step 3: Predict the charge on each pixel.}
\begin{equation}
  Q_i^{\,\text{pred}}(x_0)
    \;=\; F_i^{\,\text{pred}}(x_0) \;\cdot\; Q_{\text{total}}
  \label{eq:predicted-charge}
\end{equation}
where $Q_{\text{total}} = \sum_i Q_i^{\,\text{obs}}$ is the total
observed charge (a measured quantity, not a fit parameter).

\paragraph{Step 4: Minimise the $\chi^{2}$.}
\begin{equation}
  \chi^{2}(x_0)
    \;=\; \sum_{i}
      \frac{\bigl(Q_i^{\,\text{obs}} - Q_i^{\,\text{pred}}(x_0)\bigr)^{2}}
           {\sigma_{\text{noise},i}^{\,2}}
  \label{eq:phys-chi2}
\end{equation}
The uncertainty $\sigma_{\text{noise},i}$ is just the actual physical
noise on pixel~$i$:
\begin{equation}
  \sigma_{\text{noise},i}
    \;=\; \sqrt{\sigma_{\text{electronic}}^{\,2}
                + \bigl(\sigma_{\text{gain}} \cdot Q_i^{\,\text{obs}}\bigr)^{2}}
  \label{eq:noise}
\end{equation}
That's it.  \textbf{The only free parameter is $x_0$} (or $(x_0, y_0)$
for 2D).  The minimiser adjusts $x_0$ until the predicted charges best
match the observed charges.  The fitted $x_0$ is your reconstructed
position.

%----------------------------------------------------------------------
\subsection{Concrete Example: One Event}
\label{sec:example}
%----------------------------------------------------------------------

Suppose a hit occurs at $x_{\text{true}} = 0.12\;\text{mm}$ and the
centre pixel is at $x_{\text{px}} = 0.0\;\text{mm}$ with pitch
$p = 0.5\;\text{mm}$.  The 5 pixels in the central row are at
$x \in \{-1.0,\; -0.5,\; 0.0,\; +0.5,\; +1.0\}\;\text{mm}$.

\medskip\noindent
\textbf{Gaussian fit:} Minuit2 searches a 4D space $(A, \mu, \sigma, B)$
to find a Gaussian bell curve that best passes through the 5 charge
values.  The fitted $\mu$ (somewhere near $0.12\;\text{mm}$) is your
answer.  But $A$, $\sigma$, $B$ are nuisance parameters with no physical
meaning---they exist only because the Gaussian shape needs them to
approximate the non-Gaussian charge profile.

\medskip\noindent
\textbf{Physics-model fit:} Minuit2 searches a \textbf{1D} space
($x_0$).  For each trial value (e.g.\ $x_0 = 0.10\;\text{mm}$), it
calls your \code{calcWeight()} for all 5 pixels, normalises, multiplies
by $Q_{\text{total}}$, and compares to the observed charges.  The
$\chi^{2}$ landscape is a simple 1D curve with a clear minimum.
There is no $A$, no $\sigma$, no $B$---the physics model generates the
full predicted charge profile from $x_0$ alone.

%----------------------------------------------------------------------
\subsection{Why This Is Better}
\label{sec:why-better}
%----------------------------------------------------------------------

\begin{table}[h]
\centering
\begin{tabular}{@{}lcl@{}}
\toprule
 & \textbf{Gaussian Fit} & \textbf{Physics Model Fit} \\
\midrule
Free parameters          & 4\,(1D) or 6\,(2D)  & \textbf{1\,(1D) or 2\,(2D)} \\
NDF for 5 pixels (1D)    & 1                    & \textbf{4}                   \\
NDF for 25 pixels (2D)   & 19                   & \textbf{23}                  \\
$\sigma_{\text{err}}$ models needed
                         & 9 elaborate ones     & \textbf{Just physical noise} \\
$\chi^{2}$ physically meaningful
                         & No                   & \textbf{Yes}                 \\
Bias from model mismatch & Yes                  & \textbf{No (by construction)} \\
Fit convergence speed    & 4--6D search         & \textbf{1--2D search}         \\
\bottomrule
\end{tabular}
\end{table}

\noindent
The last row matters for performance: Minuit2 converges much faster in
1--2 dimensions than in 4--6.  Each $\chi^{2}$ evaluation calls
\code{calcWeight()} $N$ times (where $N = 5$ or $25$), which is
comparable in cost to evaluating the Gaussian $N$ times.

%----------------------------------------------------------------------
\subsection{What About $Q_{\text{total}}$?}
\label{sec:qtotal}
%----------------------------------------------------------------------

In the physics-model fit, $Q_{\text{total}}$ enters
equation~\eqref{eq:predicted-charge} as a known constant.  You have two
options:

\begin{enumerate}[nosep]
  \item \textbf{Use the observed total:}\quad
    $Q_{\text{total}} = \sum_i Q_i^{\,\text{obs}}$.  This is the
    simplest approach and works well when the neighbourhood captures most
    of the charge.  Noise on individual pixels partially cancels in the
    sum.
  \item \textbf{Fit $Q_{\text{total}}$ as a second (third) parameter:}\quad
    If charge leaks outside the neighbourhood or noise significantly
    biases the sum, you can let $Q_{\text{total}}$ float.  This adds one
    parameter (2 for 1D, 3 for 2D), still far fewer than the Gaussian.
\end{enumerate}

\noindent
For your typical $R = 2$ neighbourhood with LogA, the centre pixel
captures a large charge fraction and the sum is a good estimate, so
option~1 should suffice.

%----------------------------------------------------------------------
\subsection{Implementation Sketch}
\label{sec:implementation}
%----------------------------------------------------------------------

The implementation reuses your existing \code{ChargeSharingCore.hh}
functions.  The only new code is the $\chi^{2}$ functor and the Minuit2
wiring.

\begin{lstlisting}
#include "ChargeSharingCore.hh"
#include <Math/Minimizer.h>
#include <Math/Factory.h>
#include <Math/Functor.h>

struct PhysicsModelFitResult {
    bool converged = false;
    double x0 = NaN;       // reconstructed X
    double y0 = NaN;       // reconstructed Y
    double chi2 = NaN;
    int    ndf  = 0;
};

// pixelX[i], pixelY[i]: known pixel centre positions
// Qobs[i]:              observed charges (with noise)
// nPixels:              number of pixels (e.g. 25 for 5x5)
// pixelCenterX/Y:       centre pixel position (seed)
// pitch, padSize, d0:   detector geometry (known constants)

PhysicsModelFitResult fitPhysicsModel2D(
    const double* pixelX,  const double* pixelY,
    const double* Qobs,    int nPixels,
    double pixelCenterX,   double pixelCenterY,
    double pitch, double padSize, double d0MM,
    double sigmaElectronic, double sigmaGainFrac)
{
    // Total observed charge (known constant, not a fit parameter)
    double Qtotal = 0.0;
    for (int i = 0; i < nPixels; ++i)
        Qtotal += std::max(0.0, Qobs[i]);

    // The chi2 function: given trial (x0, y0), predict all
    // pixel charges and compare to observations.
    auto chi2Fcn = [&](const double* par) -> double {
        const double x0 = par[0];
        const double y0 = par[1];

        // Compute weights for all pixels at trial position
        double wSum = 0.0;
        double w[25];  // max neighbourhood size
        for (int i = 0; i < nPixels; ++i) {
            double dx = x0 - pixelX[i];
            double dy = y0 - pixelY[i];
            double dist = std::hypot(dx, dy);
            // YOUR EXISTING FUNCTION:
            w[i] = core::calcWeight(
                core::SignalModel::LogA,
                dist, padSize, padSize, d0MM);
            wSum += w[i];
        }

        // chi2: compare predicted vs observed
        double chi2 = 0.0;
        for (int i = 0; i < nPixels; ++i) {
            double Fpred = (wSum > 0) ? w[i] / wSum : 0.0;
            double Qpred = Fpred * Qtotal;
            double sigma = std::sqrt(
                sigmaElectronic * sigmaElectronic
                + sigmaGainFrac * sigmaGainFrac
                  * Qobs[i] * Qobs[i]);
            sigma = std::max(sigma, 1e-20);
            double pull = (Qobs[i] - Qpred) / sigma;
            chi2 += pull * pull;
        }
        return chi2;
    };

    // Set up Minuit2
    auto minimizer = std::unique_ptr<ROOT::Math::Minimizer>(
        ROOT::Math::Factory::CreateMinimizer("Minuit2", "Migrad"));
    minimizer->SetMaxFunctionCalls(200);
    minimizer->SetTolerance(1e-6);
    minimizer->SetPrintLevel(0);

    ROOT::Math::Functor functor(chi2Fcn, 2);
    minimizer->SetFunction(functor);

    // One variable per dimension, seeded at centre pixel
    minimizer->SetLimitedVariable(
        0, "x0", pixelCenterX, 0.01 * pitch,
        pixelCenterX - pitch, pixelCenterX + pitch);
    minimizer->SetLimitedVariable(
        1, "y0", pixelCenterY, 0.01 * pitch,
        pixelCenterY - pitch, pixelCenterY + pitch);

    PhysicsModelFitResult result;
    result.converged = minimizer->Minimize();
    if (result.converged) {
        result.x0   = minimizer->X()[0];
        result.y0   = minimizer->X()[1];
        result.chi2 = minimizer->MinValue();
        result.ndf  = nPixels - 2;  // 2 free params
    }
    return result;
}
\end{lstlisting}

\noindent
Note what is \emph{absent}: no amplitude~$A$, no width~$\sigma$, no
baseline~$B$, no \file{FitUncertainty.hh} models.  The physics produces
the shape; the position is the only unknown.

%----------------------------------------------------------------------
\subsection{Relationship to the Gaussian Fit}
\label{sec:relationship}
%----------------------------------------------------------------------

The Gaussian fit is doing the same thing conceptually---minimising a
$\chi^{2}$ between a model and data---but with the wrong model.  The
table below shows the correspondence:

\medskip
\begin{center}
\begin{tabular}{@{}lll@{}}
\toprule
\textbf{Concept} & \textbf{Gaussian fit} & \textbf{Physics-model fit} \\
\midrule
Fit function
  & $A\, e^{-(x-\mu)^{2}/2\sigma^{2}} + B$
  & $F_i(x_0) \cdot Q_{\text{total}}$ \\[2pt]
Position parameter
  & $\mu$
  & $x_0$ \\[2pt]
Nuisance parameters
  & $A,\, \sigma,\, B$
  & \emph{none} \\[2pt]
Model shape
  & Gaussian (wrong)
  & LogA/LinA (correct) \\[2pt]
Per-pixel error
  & 9 tuneable models
  & Physical noise \\[2pt]
Minimiser
  & Minuit2 (reuse as-is)
  & Minuit2 (reuse as-is) \\
\bottomrule
\end{tabular}
\end{center}

\medskip\noindent
You keep the same minimiser, the same convergence-check logic, the same
centroid fallback.  The only change is what the $\chi^{2}$ function
computes inside.

%----------------------------------------------------------------------
\section{Validation: Tornago et al.\ (arXiv:2007.09528)}
\label{sec:tornago}
%----------------------------------------------------------------------

The physics-model fit described in Section~\ref{sec:physfit} is not a
new idea---it is exactly the reconstruction strategy derived and
validated in the foundational AC-LGAD paper by Tornago et al.~\cite{tornago2021}.

\subsection{What the Paper Does}

Section~3.3 of Tornago et al.\ defines position reconstruction as
follows.  For each event, the measured charge fractions are:
\begin{equation}
  f_i^{\,\text{meas}} \;=\; \frac{A_i}{A_{\text{tot}}}
  \label{eq:tornago-fmeas}
\end{equation}
where $A_i$ is the amplitude (charge) on pad~$i$ and
$A_{\text{tot}} = \sum_j A_j$.  The model-predicted fractions for a
trial position $(x, y)$ are computed from the LogA or LinA model
(Eqs.~4 and~6 in the paper).  The reconstructed position is the one
that minimises the $\chi^{2}$ (their Eq.~8):
\begin{equation}
  \chi^{2}(x, y)
    \;=\; \sum_{i}
      \left[
        \left(\frac{A_i}{A_{\text{tot}}}\right)_{\!\text{meas}}
        \;-\;
        \left(\frac{A_i}{A_{\text{tot}}}\right)_{\!\text{calc}}
      \right]^{2}
  \label{eq:tornago-chi2}
\end{equation}

\noindent
This is the same objective as equation~\eqref{eq:phys-chi2} in
Section~\ref{sec:phys-fit}: compare measured charge fractions against
fractions predicted by the physics model, and find the position that
best explains the data.

\subsection{Difference: Grid Search vs.\ Continuous Minimisation}

Tornago et al.\ perform the minimisation as a \textbf{grid search}:
the detector surface is divided into discrete $(x, y)$ bins, the
predicted fractions are precomputed for each bin, and the bin with the
smallest $\chi^{2}$ is selected.  The paper notes that ``the accuracy
can be increased by performing a local interpolation around the
minimum.''

The approach in Section~\ref{sec:phys-fit} replaces the grid search
with \textbf{continuous Minuit2 minimisation}, which has several
advantages:

\begin{itemize}[nosep]
  \item \textbf{Sub-bin precision.}\quad
    The reconstructed position is not quantised to a grid.  Minuit2
    naturally finds the exact minimum of the $\chi^{2}$ surface.
  \item \textbf{No precomputation.}\quad
    There is no need to build and store a lookup table of predicted
    fractions for every bin.  The forward model is evaluated on-the-fly
    during minimisation.
  \item \textbf{Proper uncertainties.}\quad
    Per-pixel noise enters the denominator of the $\chi^{2}$
    (equation~\eqref{eq:phys-chi2}), giving the minimiser correct
    weighting.  The paper's Eq.~8 uses unweighted residuals.
  \item \textbf{Fit quality metric.}\quad
    The minimised $\chi^{2}/\text{NDF}$ is a meaningful goodness-of-fit
    measure, enabling quality cuts.
\end{itemize}

\noindent
Both approaches are valid.  The grid search is simpler to implement and
avoids any risk of minimiser non-convergence; the continuous fit gives
better resolution and richer diagnostics.

\subsection{Connection to the Codebase}

The functions used by Tornago et al.\ to compute the predicted fractions
are the \emph{same} functions already implemented in
\file{ChargeSharingCore.hh}:

\medskip
\begin{center}
\begin{tabular}{@{}lll@{}}
\toprule
\textbf{Tornago et al.} & \textbf{Your code} & \textbf{Purpose} \\
\midrule
Eq.~4 (LogA weight)
  & \code{calcWeightLogA()}
  & $F_i \propto \alpha_i / \ln(d_i/d_0)$ \\[2pt]
Eq.~6 (LinA weight)
  & \code{calcWeightLinA()}
  & $F_i \propto (1 - \beta\, d_i)\,\alpha_i$ \\[2pt]
Fraction normalisation
  & \code{calculateNeighborhood()}
  & $\sum_i F_i = 1$ \\[2pt]
Pad view angle $\alpha_i$
  & \code{calcPadViewAngle()}
  & Geometric acceptance \\
\bottomrule
\end{tabular}
\end{center}

\medskip\noindent
In other words, the codebase already contains a complete implementation
of the Tornago reconstruction model on the \emph{simulation} side.
The only missing piece is calling these same functions on the
\emph{reconstruction} side---inside a Minuit2 $\chi^{2}$ functor---rather
than fitting a Gaussian.

\begin{thebibliography}{1}
\bibitem{tornago2021}
  N.~Cartiglia, M.~Tornago, et al.,
  ``Resistive AC-Coupled Silicon Detectors: principles of operation,
    flavors, performance,''
  \textit{Nucl.\ Instrum.\ Meth.\ A}, vol.~1003, p.~165311, 2021.
  arXiv:2007.09528.
\end{thebibliography}

%----------------------------------------------------------------------
\section{Secondary Issues}
%----------------------------------------------------------------------

\paragraph{1.\ \file{FitUncertainty.hh} is over-engineered.}
Nine $\sigma$ models (\code{SigmaLinear}, \code{SigmaPower},
\code{SigmaPowerInverse}, \code{SigmaQuadratic},
\code{SigmaExpSaturating}, \code{SigmaPiecewiseCoreTail},
\code{SigmaLogistic}, \code{QnQiScaled}, \code{UniformPercentOfMax})
with ${\sim}\,15$ tuneable parameters.  This is a hyperparameter search
to compensate for Gaussian model mismatch.  With the correct fit
function, you need one model: physical noise.

\paragraph{2.\ No fit quality cut.}
A fit that ``converges'' but has $\sigma$ pegged at a boundary, or
$\chi^{2}/\text{NDF} \gg 1$, or $A < 0$, is not reliable.  The plugin's
centroid fallback triggers on non-convergence only.  Add rejection
criteria:
\begin{itemize}[nosep]
  \item $\sigma$ at boundary (pegged) $\;\to\;$ reject.
  \item $\chi^{2}/\text{NDF} > \text{threshold}$ $\;\to\;$ reject.
  \item $A$ or $B$ at boundary $\;\to\;$ suspect.
\end{itemize}

\paragraph{3.\ 2D Gaussian lacks a rotation angle.}
Your 2D Gaussian assumes axis-aligned elliptical contours (separate
$\sigma_x$, $\sigma_y$).  For hits not on the pixel centre axis, the
charge distribution has a rotated ellipse shape.  Adding a rotation
parameter~$\theta$ (7 parameters total) would help, but this is minor
compared to the model mismatch issue.

\paragraph{4.\ Standalone fits are batch post-processing, not inline.}
\file{GaussianFitter.cc} opens the ROOT file after simulation and adds
branches.  This means during event processing, reconstructed positions
aren't available.  The plugin does it correctly (inline).
\file{PROJECT\_REPORT.md} already flags this.

%----------------------------------------------------------------------
\section{Do You Need to Improve?}
%----------------------------------------------------------------------

\textbf{For simulation studies / algorithm R\&D:}\quad
The Gaussian fits are \emph{adequate}.  They give reasonable position
resolution for benchmarking charge sharing models and comparing LogA
vs.\ LinA.  The centroid fallback ensures robustness.

\medskip\noindent
\textbf{For ePIC production (integration into EICrecon):}\quad
Yes, improving matters.  Recent results from the community:

\begin{itemize}[nosep]
  \item \textbf{Siviero et al.\ (2024):} Neural network on FBK RSD2
    ($450\;\mu\text{m}$ pitch) achieved ${\sim}65\;\mu\text{m}$ resolution
    at DESY test beam---${\sim}50\%$ better than analytical methods.
  \item \textbf{Sunnarborg, CPAD 2025:} DNN on BNL AC-LGADs
    ($500\;\mu\text{m}$ pitch) achieved ${\sim}20\;\mu\text{m}$ MSE
    vs.\ ${\sim}25\;\mu\text{m}$ for analytical matrix inversion.
  \item Both groups found that ML on full waveforms beats analytical
    amplitude-only methods.
\end{itemize}

\noindent
The physics-model fit described in Section~\ref{sec:physfit} would be a
natural middle ground: better than a Gaussian, simpler than training an
NN, and reusing your existing code.

%----------------------------------------------------------------------
\section{Summary of Priorities}
%----------------------------------------------------------------------

\begin{enumerate}[nosep]
  \item \textbf{High value, moderate effort:}\quad
    Replace the Gaussian with your own physics model as the fit function.
    You already have \code{calcWeight()}---wrap it in a Minuit2
    \code{FCNBase}.  This eliminates \file{FitUncertainty} complexity and
    improves accuracy.

  \item \textbf{Quick win:}\quad
    Add fit quality cuts (boundary-pegged parameters,
    $\chi^{2}/\text{NDF}$ threshold) to the reconstruction fallback
    logic.

  \item \textbf{Low priority:}\quad
    The Gaussian fit is fine to keep as a \emph{secondary/comparison}
    method---it's useful for visualisation and diagnostics.  But it
    shouldn't be the primary reconstruction.

  \item \textbf{Future:}\quad
    For ePIC production, consider that ML approaches are achieving the
    best resolution in the community.  Your simulation is ideal training
    data for a lightweight neural network.
\end{enumerate}

\end{document}
